\DocumentMetadata{
  pdfversion=2.0,
  pdfstandard=ua-2,
  lang=en-US,
  tagging=on,
  tagging-setup={math/setup=mathml-SE,tabsorder=structure}
}
\documentclass[11pt,letterpaper]{article}
\usepackage[margin=0.68in,includefoot]{geometry}
\usepackage{newcomputermodern}
\usepackage{amsmath,amscd,mathtools}
\usepackage{tikz,adjustbox}
\providecommand{\square}{\mdlgwhtsquare}
\usepackage[hidelinks]{hyperref}
\hypersetup{
  pdftitle={Numerical Analysis PhD Exam, January 2010},
  pdfauthor={University of Florida Department of Mathematics},
  pdfsubject={Graduate mathematics examination},
  pdfdisplaydoctitle=true,
  bookmarksopen=true
}
\setcounter{secnumdepth}{0}
\setlength{\parindent}{0pt}
\setlength{\parskip}{0.28em}
\setlength{\emergencystretch}{3em}
\setlength{\leftmargini}{2.15em}
\setlength{\leftmarginii}{2.65em}
\setlength{\leftmarginiii}{3em}
\renewcommand{\labelenumi}{\textbf{\theenumi.}}
\pagestyle{empty}
\raggedbottom
\allowdisplaybreaks
\newenvironment{examproblems}[1]{%
  \begin{enumerate}%
  \setcounter{enumi}{#1}%
  \setlength{\topsep}{0.35em}%
  \setlength{\partopsep}{0pt}%
  \setlength{\itemsep}{0.48em}%
  \setlength{\parsep}{0.18em}%
}{\end{enumerate}}
\newenvironment{examsubparts}{%
  \begin{enumerate}%
  \setlength{\topsep}{0.18em}%
  \setlength{\partopsep}{0pt}%
  \setlength{\itemsep}{0.16em}%
  \setlength{\parsep}{0.1em}%
}{\end{enumerate}}
\newenvironment{examinstructions}{%
  \par\begingroup\itshape
}{%
  \par\endgroup\vspace{0.18em}
}
\makeatletter
\renewcommand\section{\@startsection{section}{1}{0pt}%
  {-1.25ex plus -.25ex minus -.1ex}%
  {0.55ex plus .1ex}%
  {\normalfont\large\bfseries}}
\renewcommand{\@maketitle}{%
  \newpage\null\vskip -1em%
  {\footnotesize\bfseries
    \href{https://www.ufl.edu/}{University of Florida}, \href{https://math.ufl.edu/}{Department of Mathematics}\par}%
  \vskip -0.5em%
  \tagstructbegin{tag=Title}%
    {\LARGE\bfseries\href{https://gma.math.ufl.edu/exams/numerical-analysis-phd/}{Numerical Analysis PhD Exam}, January 2010\par}%
  \tagstructend%
  \vskip 0.25em%
  \tagmcbegin{artifact}%
    \hrule height 1pt%
  \tagmcend%
  \vskip 1em%
}
\makeatother
\title{Numerical Analysis PhD Exam, January 2010}
\author{}
\date{}
\begin{document}
\maketitle
\thispagestyle{empty}
\begin{examinstructions}
Answer any 8 questions.
\end{examinstructions}
\begin{examproblems}{0}
\item
Prove that a projector is normal if and only if it is self-adjoint.
\item
Let \(A \in \mathbb{R}^{N\times N}\) and \(b \in \mathbb{R}^N\). Consider the following iteration for solving \(Ax=b\), that computes \(x_{n+1}\), given \(x_n \in \mathbb{R}^N\), as follows (in~\(m\) intermediate steps): Setting \(x_{n+1}^{(0)}=x_n\), for \(\ell=1,2,\ldots,m\), compute \(x_{n+1}^{(\ell)}=x_{n+1}^{(\ell-1)}+\tau_\ell(b-Ax_{n+1}^{(\ell-1)})\). Then, define \(x_{n+1}=x_{n+1}^{(m)}\). There is a linear operator~\(E\) such that \(x_{n+1}-x=E(x_n-x)\). Give a formula for~\(E\). Suppose~\(A\) is Hermitian and positive definite with spectral condition number~\(\kappa\). Prove that there are real values of the~\(m\) parameters \(\tau_\ell\) such that 
\[
\rho(E)\leq 2\left(\frac{\sqrt{\kappa}-1}{\sqrt{\kappa}+1}\right)^m.
\]
\item
Suppose~\(A\) and~\(B\) are square matrices such that \(AB\) is normal. Prove that \(\lVert AB\rVert_2\leq \lVert BA\rVert\). (We use \(\lVert\cdot\rVert_2\) to denote the spectral norm and \(\lVert\cdot\rVert\) to denote any induced matrix norm.)
\item
Let \(A\in\mathbb{C}^{m\times m}\), and \(a_j\) be its~\(j\)-th column. Prove that 
\[
|\det A|\leq \prod_{j=1}^m\lVert a_j\rVert_2.
\]
\item
Suppose~\(A\) is a Hermitian positive definite matrix split into \(A=C+C^*+D\) where~\(D\) is also Hermitian positive definite. Show that \(B=C+\omega^{-1}D\) is invertible. Consider the iteration \(x_{n+1}=x_n+B^{-1}(b-Ax_n)\), with any initial iterate \(x_0\). Prove that \(x_n\) converges to \(x=A^{-1}b\) whenever \(0<\omega<2\).
\item
Let \(x_0,x_1,\ldots,x_n\) be distinct points in a finite interval \([a,b]\) and \(f\in C^1[a,b]\). Show that for any given \(\varepsilon>0\) there exists a polynomial~\(p\) such that \(\lVert f-p\rVert_\infty<\varepsilon\) and \(p(x_i)=f(x_i)\) for all \(i=0,1,\ldots,n\) (where \(\lVert\cdot\rVert_\infty\) denotes the \(L^\infty(a,b)\)-norm).
\item
Let \(p>0\) and \(x=\sqrt{p+\sqrt{p+\sqrt{p+\cdots}}}\), where all the square roots are positive. Design a fixed point iteration \(x_{n+1}=F(x_n)\) that converges to~\(x\). You should write down a specific~\(F\). Obtain a sufficient condition on the initial iterates for (global) convergence of the iteration.
\item
Let~\(n\) be a positive integer, \(h=1/n\), and consider the grid of points \((ih,jh)\) for \(i,j=0,1,\ldots,n\). Let~\(A\) be the finite difference operator with the “5-point stencil” discretizing the Laplace operator \(-\Delta\), with zero Dirichlet boundary conditions on this grid. Describe it. Prove that the spectrum of~\(A\) consists of the numbers 
\[
\lambda_{lm}=4h^{-2}\bigl(\sin^2(l\pi h/2)+\sin^2(m\pi h/2)\bigr),
\]
 for all \(l,m=1,\ldots,n-1\).
\item
Let \(x_m\) and \(x_{m+1}\) be two successive (complex) iterates when Newton’s method is applied to a polynomial \(p(z)\) of degree~\(n\). Prove that there is a zero of \(p(z)\) in the disk \(\{z\in\mathbb{C}:|z-x_m|\leq n|x_{m+1}-x_m|\}\).
\item
Let \(w(x)\) be a positive integrable function on \([a,b]\). Consider the quadrature 
\[
\int_a^b f(x)w(x)\,dx\approx\sum_{k=0}^n\bigl(A_kf(x_k)+B_kf'(x_k)+C_kf''(x_k)\bigr)
\]
 for any~\(f\) with continuous first and second derivatives (\(f'\) and \(f''\), respectively). Give conditions on \(A_k,B_k,C_k\) and \(x_k\) so that the quadrature has precision \(4n+3\).
\end{examproblems}
\end{document}
